%!TEX root = ../main.tex

\section{Устойчивость положений равновесия системы}

\begin{frame}{Одномерная задача}
\vspace{-5mm}
Работа \cite{ponomarev_stability}
\vspace{1mm}
\begin{itemize}
	\item $\phi(\vx, 0) = \phi_0(\vx) = \phi_0(x)$, $\epsilon_0(\vx) = \epsilon_0(x)$ не зависят от $y$ и $z$
	\item $\Phi|_{y = 0} = \Phi^- \in \Real$, $\Phi|_{y = h} = \Phi^+ \in \Real$
\end{itemize}
\vspace{1mm}
Функция электрического потенциала известна аналитически:
\begin{columns}
\column{0.65\textwidth}
\vspace{-4mm}
\[
	\Phi(\vx, t) = \Phi^- \hm + \cfrac{y}{h}(\Phi^+ - \Phi^-), \quad E = \| \nabla \Phi \| = \Const.
\]
\raggedleft
\begin{minipage}{0.98\textwidth}
\vspace{-2mm}
\begin{block}{Уравнение фазового поля $\phi$}
	\[
		\cfrac{1}{m} \partt{\phi} = \half E^2 \epsilon'(\phi) + \cfrac{\Gamma}{l^2} f'(\phi) + \half \Gamma \cfrac{\partial^2 \phi}{\partial x^2}
	\]
	\vspace{-4mm}
\end{block}
\end{minipage}
\column{0.35\textwidth}
\begin{figure}
	\vspace*{-18mm}
	\hspace*{2mm}
	\includegraphics[width=0.93\textwidth]{figures/one_dim_problem.jpg}
\end{figure}
\end{columns}
\end{frame}


\begin{frame}{Анализ положений равновесия}
\begin{itemize}
	\item Канал пробоя может развиваться из малых возмущений неповрежденной среды.
	\item Рассмотрим положения равновесия вида $\phi(x, t) \equiv C$. \\
	Положению равновесия соответствует ноль $C$ функции
\end{itemize}
\[
	\chi(\phi) = \half E^2 \epsilon'(\phi) + \cfrac{\Gamma}{l^2} f'(\phi).
\]
\vspace{-0.5cm}
\begin{itemize}
	\item Исследуем положения равновесия на устойчивость спектральным методом: к $\phi \equiv C$ прибавим возмущение $\delta \phi = e^{\alpha t} \sin(\omega x)$, линеаризуем уравнение на $\delta \phi$.
\end{itemize}
\begin{center}
\begin{tabular}{c c c}
	$\chi(\phi)$ возрастает в $C$ & $\Longrightarrow$ & равновесие неустойчиво; \\
	$\chi(\phi)$ убывает в $C$ & $\Longrightarrow$ & равновесие устойчиво.
\end{tabular}
\end{center}
\end{frame}


{\nologo
\begin{frame}{Анализ положений равновесия}
\vspace{-8mm}
\begin{columns}
\column{0.3\textwidth}
	\begin{center}
		<<Слабое>> напряжение
	\end{center}
	\vspace{-1mm}
	\[
		0 \leqslant \cfrac{E^2 l^2 \epsilon_0}{2 \Gamma} < \delta^2
	\]
	\vspace{-6mm}
	\begin{figure}
		\includegraphics[width=\textwidth]{figures/equilibriums_case_1.png}
	\end{figure}
\column{0.3\textwidth}
	\begin{center}
		<<Среднее>> напряжение
	\end{center}
	\vspace{-1mm}
	\[
		\delta^2 < \cfrac{E^2 l^2 \epsilon_0}{2 \Gamma} < (1 + \delta)^2
	\]
	\vspace{-6mm}
	\begin{figure}
		\includegraphics[width=\textwidth]{figures/equilibriums_case_2.png}
	\end{figure}
\column{0.3\textwidth}
	\begin{center}
		<<Сильное>> напряжение
	\end{center}
	\vspace{-1mm}
	\[
		(1 + \delta)^2 < \cfrac{E^2 l^2 \epsilon_0}{2 \Gamma}
	\]
	\vspace{-6mm}
	\begin{figure}
		\includegraphics[width=\textwidth]{figures/equilibriums_case_3.png}
	\end{figure}
\end{columns}
\vspace{1mm}
\begin{columns}
\column{0.3\textwidth}
	\hspace{5mm}
	$\phi \equiv 0$ неустойчивое \\
	\hspace{5mm}
	$\phi \equiv 1$ устойчивое
\column{0.3\textwidth}
	\hspace{5mm}
	$\phi \equiv 0$ устойчивое \\
	\hspace{5mm}
	$\phi \equiv С_3$ неустойчивое \\
	\hspace{5mm}
	$\phi \equiv 1$ устойчивое
\column{0.3\textwidth}
	\hspace{5mm}
	$\phi \equiv 0$ устойчивое \\
	\hspace{5mm}
	$\phi \equiv 1$ неустойчивое
\end{columns}
\end{frame}
}


\begin{frame}{Разностная схема}
\begin{block}{Разностная задача}
	$$\cfrac{1}{m} \cfrac{\phi_i^{j + 1} - \phi_i^j}{\tau} = \half E_\phi^2 \epsilon'(\phi_i^j) + \cfrac{\Gamma}{l^2} f'(\phi_i^j) + \cfrac{\Gamma}{2} \cfrac{\phi_{i + 1}^j - 2 \phi_i^j + \phi_{i - 1}^j}{h^2}$$
	$$\phi_i^0 = \phi_0(ih); \quad \phi_0^j = \phi_l(j \tau); \quad \phi_n^j = \phi_r(j \tau)$$
	Сетка регулярная; $\tau$ -- шаг по времени, $h$ -- шаг по пространству.
\end{block}
\vspace{1mm}
Явная разностная схема первого порядка по времени, второго -- по пространству.
\end{frame}


\begin{frame}{Оценка устойчивости}
\begin{center}
\begin{minipage}{0.95\textwidth}
\begin{block}{Условие устойчивости}
	$$\tau \leqslant \cfrac{1}{2 m} \left( \cfrac{E^2 \epsilon_0}{\delta^{5/3}} + \cfrac{\Gamma}{h^2} \right)^{-1}$$
\end{block}
\begin{block}{Упрощенное условие устойчивости}
	$$\tau \leqslant \cfrac{1}{4 m} \min \left(\cfrac{\delta^{5/3}}{E^2 \epsilon_0}, \; \cfrac{h^2}{\Gamma} \right)$$
\end{block}
\end{minipage}
\end{center}
Оценка получена применением спектрального признака устойчивости \\ в окрестности $\phi = 0$.
\end{frame}


\begin{frame}{Вычисления: типичное решение}
\vspace{-0.4cm}
\begin{columns}
\column{0.88\textwidth}
\begin{figure}
	\includegraphics[width=\textwidth]{figures/typical_solution.png}
\end{figure}
\column{0.12\textwidth}
\hfill \\
\vspace{3.5cm}
\hspace{-2.5cm}
Узлов по измерениям: \\
\hspace{-2.5cm}
$N_x = 10^3$, $N_t = 10^5$
\end{columns}
\end{frame}


\begin{frame}{Вычисления: проверка устойчивости}
\vspace{-0.4cm}
\begin{columns}
\column{0.7\textwidth}
\begin{figure}
	\includegraphics[width=\textwidth]{figures/stability_bounds.png}
\end{figure}
\column{0.3\textwidth}
\hspace{1mm} Условие устойчивости:
\[
	\tau \leqslant \cfrac{1}{2m} \left( \cfrac{E^2 \epsilon_0}{\delta^{5/3}} + \cfrac{\Gamma}{h^2} \right)^{-1}
\]
\end{columns}
\end{frame}